\section{Integrative synthesis: from perception to response persistence}
\label{sec:synthesis}

% Receiver-centered model and critical distinctions.
% See docs/OUTLINE.md — Integrative synthesis; Proposed figures and table.
% Draft 1. Style model: Blumstein (agents/STYLE.md). Figures are placeholders (docs/FIGURE_PROMPTS.md).

The questions reviewed above converge on one observation: the same acoustic event can elicit approach, vigilance, freezing, an endocrine response without overt behavior, or nothing at all, depending on who hears it, what that receiver has experienced, and when the sound occurs. Laboratory and field studies have documented this variation largely in isolation, using different species, stimuli, and endpoints. We propose a receiver-centered framework that places these results on a common footing (\figref{fig:framework}). Its organizing currency is functional: when responding is cheap relative to missing a real threat, selection should tolerate many unnecessary responses \citep{Nesse2005}, and the relevant question becomes when receivers revise the meaning they assign to a sound rather than whether they respond.

\subsection{Four levels of receiver processing}
\label{sec:synthesis-levels}

The framework separates acoustic detection, biological interpretation, action selection, and updating. The separation is warranted because each pair of adjacent levels has been experimentally dissociated, and a measurement taken at one level therefore cannot stand in for another.

At the level of detection, receivers respond to acoustic features that need not carry species-specific meaning. In naïve rats, amygdala neurons fired differently to 22-kHz and 50-kHz calls, yet frequency-matched tones reproduced most of this difference \citep{ParsanaEtAl2012a}. Behaviorally, sine-wave approximations of 50-kHz calls that lacked natural frequency and amplitude modulation were sufficient to elicit approach in adult rats \citep{WohrSchwarting2007}. A differential response to two calls thus establishes discrimination of their acoustic properties, not that the receiver assigns them different biological meaning.

Interpretation draws on both intrinsic bias and experience, and their relative contributions differ among systems. Selective freezing to 22-kHz calls depends on prior aversive experience \citep{KimEtAl2010,EndresEtAl2007} and persists when the receiver was prevented from vocalizing during that experience \citep{CalubEtAl2018}. Maternal experience broadens behavioral and cortical tuning to the temporal structure of pup calls on top of an intrinsic bias toward species-typical intervals \citep{SchiavoEtAl2020}. For avian predator calls, results from predator-naïve laboratory rodents conflict: DBA/2 mice responded defensively to owl calls \citep{HendrieNeill1991}, whereas C57BL/6J mice, rats, and gerbils showed no selective response \citep{KindermannEtAl2009}, so strain as well as paradigm may contribute.
Field studies add dimensions that laboratory designs rarely manipulate. Marmot alarm calls encode urgency through call rate rather than predator identity \citep{BlumsteinArmitage1997}, and several rodents discriminate individual callers \citep{Hare1998,BlumsteinDaniel2004,VielleEtAl2021,WarrenEtAl2025}.

Interpretation, in turn, does not fix the action. Owl-call playback raised glucocorticoids in both voles and spiny mice, but only voles froze or fled \citep{EilamEtAl1999}. In laboratory mice, the same synthetic frequency sweep produced flight or freezing depending on whether the animal was in its home cage and whether it had previously been shocked \citep{MongeauEtAl2003}. These outcomes match the economics of escape, in which the response chosen depends on the cost of fleeing relative to the cost of staying \citep{YdenbergDill1986}. Field studies supply the costs that laboratory assays omit: alarm calls delayed re-emergence and reduced food intake in chipmunks \citep{BaackSwitzer2000}, and predator-call playback suppressed and delayed foraging in gray squirrels without changing distance to refuge \citep{JayneEtAl2015}. The absence of freezing is therefore weak evidence that a sound was not assessed.

Neural measures do not map neatly onto a single level. In head-fixed mice, the same aversive noise lowered dopamine in the nucleus accumbens while raising it in the tail of the striatum, and the accumbens decrease grew with noise intensity \citep{FuruyamaEtAl2026}. In freely moving rats, by contrast, accumbens dopamine did not scale with noise intensity even though avoidance did \citep{GoedhoopEtAl2022}. Whether the difference reflects species, recording method, or the opportunity to escape is unresolved, but neither result supports reading accumbens dopamine as a scalar measure of threat.
In Richardson's ground squirrels, alarm chirps reduced rather than increased c-Fos labeling in the medial amygdala and superior colliculus \citep{FreemanEtAl2020}. Neural responses are best treated as evidence about which computations a receiver performs, to be interpreted alongside the behavioral level at which the corresponding decision is expressed.

\begin{figure}[tbp]
  \centering
  \includegraphics[width=\textwidth]{fig1_framework_placeholder.png}
  \caption{Receiver-centered framework. (A) Experiments probe four levels that can dissociate: acoustic detection, biological interpretation, action selection, and updating. Prior experience, ecological context, and caller identity and reliability shape interpretation; the relative costs of missing a threat and of a false alarm shape action selection. Double bars mark documented dissociations: amygdala neurons distinguish call types largely by frequency and without freezing \citep{ParsanaEtAl2012a}, and spiny mice show endocrine but not behavioral responses to owl calls \citep{EilamEtAl1999}. The dashed link (temporal expectation) has not been tested with natural rodent calls. (B) The same class of sound carries different meaning for receivers with different histories \citep{HendrieNeill1991,KindermannEtAl2009,KimEtAl2010,SchiavoEtAl2020,HareAtkins2001,BlumsteinEtAl2004}.}
 
  \label{fig:framework}
\end{figure}

\subsection{Updating: distinguishing processes that reduce responding}
\label{sec:synthesis-updating}

A declining response across repeated playbacks is among the most common results in this literature, and also among the least diagnostic. At least four processes produce it: sensory habituation, learning that a signal is unreliable, a fall in the motivational significance of the stimulus, and growing expectation of when the next sound will occur. Formal habituation criteria, including stimulus specificity, spontaneous recovery, and dishabituation, exclude fatigue and sensory adaptation \citep{RankinEtAl2009} but do not by themselves separate the remaining three. Each, however, predicts a different pattern of decline and recovery (\figref{fig:temporal}).

If ecological significance slows habituation or speeds recovery, natural calls should retain or regain responsiveness relative to acoustically matched controls. The most direct test found exactly this asymmetry for recovery: after rests of one or seven days, rat orienting recovered more to a loud natural distress squeal than to its mimic \citep{EvansHammond1983a}.
Significance does not, however, prevent habituation. Rat approach to prosocial 50-kHz calls habituated even when playback reliably predicted social interaction \citep{BerzEtAl2022}, and both approach and accumbens dopamine responses to these calls declined rapidly with repetition \citep{WilluhnEtAl2014}. The comparison that would establish whether significance changes the rate of habituation, natural calls versus matched synthetic controls under identical schedules, remains largely untested outside that single study.

Reliability learning makes a different prediction: the decline should be specific to the source and should reverse when the source changes. In Richardson's ground squirrels, vigilance fell across repeated calls from one individual and was restored by a call from a different individual but not by a new call from the same individual \citep{Hare1998}, and juveniles discounted callers whose alarms had repeatedly gone unaccompanied by a predator model \citep{HareAtkins2001}.
Yellow-bellied marmots, by contrast, responded more to experimentally unreliable callers \citep{BlumsteinEtAl2004}, and caller familiarity alone had no detectable effect \citep{BlumsteinDaniel2004}. Discounting and increased scrutiny are both plausible solutions to unreliable information; which is favored should depend on the cost of independently checking for danger relative to the cost of ignoring a true alarm, a prediction that has not been tested.

Temporal expectation predicts that the schedule of exposure, in addition to the number of exposures, shapes the decline. At a matched mean interval, regularly timed startle stimuli produced more habituation than irregularly timed ones \citep{Davis1970}, and temporally jittered tone pulses prevented rapid habituation of basolateral amygdala responses \citep{HerryEtAl2007}. Both studies used artificial sounds, and neither compared distributions of intervals, so whether these effects extend to natural calls is unknown. At a longer timescale, the risk allocation hypothesis predicts that vigilance per alarm should fall when high-risk periods are frequent and rise when they are rare \citep{LimaBednekoff1999}; it has not been tested with acoustic cues in rodents.

\begin{figure}[tbp]
  \centering
  \includegraphics[width=\textwidth]{fig2_temporal_placeholder.png}
  \caption{Hypothesized temporal dynamics of responses to repeated sounds; curves are schematic, not data. (A) Irregular timing slows within-session decline for artificial sounds \citep{Davis1970,HerryEtAl2007}; whether natural calls and matched controls diverge under each schedule is untested. (B) After rest, orienting recovered more to a loud natural distress call than to its acoustic mimic \citep{EvansHammond1983a}. (C) Receivers track caller reliability, but in opposite directions: juvenile ground squirrels discount unreliable callers \citep{HareAtkins2001}, whereas marmots respond more to them \citep{BlumsteinEtAl2004}; a change of caller restores responding \citep{Hare1998}. (D) Risk allocation predicts less vigilance per alarm when high-risk periods are frequent \citep{LimaBednekoff1999}; no rodent acoustic test exists.}
 
  \label{fig:temporal}
\end{figure}

These accounts are not mutually exclusive, and a single playback series may engage all of them. The framework's value lies in specifying what each predicts: stimulus-specific decline with recovery after rest (habituation), source-specific decline with renewal when the source changes (reliability), and schedule-dependent decline at a matched number of exposures (temporal expectation). Each exposure may therefore update three quantities, the expected timing of the next sound, its expected consequence, and the reliability of its source. This is an inference across studies rather than a demonstrated result, and \secref{sec:conclusion} develops the experiments that could test it.

% TODO: critical-distinctions table from docs/OUTLINE.md (Integrative synthesis) not yet drafted.
